\section{Wstęp}
\subsection{Uzasadnienie wyboru tematu}
% [PRZECZYTANE] [PRZECZYTANE] [PRZECZYTANE] [PRZECZYTANE] [PRZECZYTANE]
Technologie Speech-to-Text (STT) coraz częściej stają się elementem interfejsów użytkownika w aplikacjach webowych, mobilnych i asystenckich. W praktyce oznacza to, że dokładność transkrypcji, czas odpowiedzi oraz odporność na zakłócenia mają bezpośredni wpływ na komfort użytkownika i przydatność całego systemu. Szczególnie istotne jest to w przypadku języka polskiego, dla którego dostępność dopracowanych rozwiązań jest nadal mniejsza niż dla języka angielskiego, a wyniki badań wskazują na duże zróżnicowanie jakości pomiędzy poszczególnymi modelami i usługami \cite{zielonka2023survey, pondel2024e2e,ferraro2023benchmarking}.

Temat pracy jest uzasadniony z perspektywy praktycznej, naukowej i metodologicznej. Z perspektywy praktycznej rośnie potrzeba projektowania aplikacji webowych działających w czasie rzeczywistym, w których rozpoznawanie mowy może wspierać na przykład dyktowanie, obsługę systemów bez użycia klawiatury oraz szybkie prototypowanie interfejsów głosowych. Z perspektywy naukowej nadal widoczna jest luka badawcza dotycząca porównań rozwiązań STT dla języka polskiego w warunkach zbliżonych do rzeczywistych aplikacji webowych, z uwzględnieniem zarówno modeli offline, jak i usług chmurowych oraz natywnych mechanizmów przeglądarkowych. Z perspektywy metodologicznej temat pozwala połączyć implementację systemu eksperymentalnego z obiektywną oceną jakości i wydajności, co czyni go dobrym materiałem do pracy inżyniersko-badawczej \cite{chlewicki2024potencjal,pudo2024optimizing}.

W literaturze coraz częściej pojawiają się analizy porównawcze rozwiązań STT, jednak wiele z nich koncentruje się na innych językach, na modelach uruchamianych poza przeglądarką albo na jednym wybranym aspekcie jakości. W przypadku języka polskiego szczególnie wartościowe jest zestawienie rozwiązań działających lokalnie i w chmurze w jednym, powtarzalnym środowisku testowym, ponieważ dopiero takie podejście pozwala realnie ocenić ich przydatność do zastosowań webowych.

\subsection{Cel i zakres pracy}
% [PRZECZYTANE] [PRZECZYTANE] [PRZECZYTANE] [PRZECZYTANE] [PRZECZYTANE]
Celem głównym pracy jest przeprowadzenie porównawczej analizy technologii Speech-to-Text pod kątem ich użyteczności w webowych aplikacjach dla języka polskiego. Analiza obejmuje zarówno jakość transkrypcji, jak i zachowanie systemów w warunkach zbliżonych do praktycznego użycia, czyli przy zmiennym poziomie zakłóceń, różnym tempie mówienia oraz odmiennych charakterystykach głosów.

Cele szczegółowe obejmują ocenę dokładności transkrypcji z wykorzystaniem metryk WER oraz CER, pomiar opóźnień przetwarzania w ujęciu RTT, badanie odporności rozwiązań na szum i zróżnicowane warunki wejściowe, a także przygotowanie webowej aplikacji eksperymentalnej umożliwiającej wykonywanie porównywalnych testów. W ramach pracy zakłada się porównanie kilku klas rozwiązań, w tym mechanizmów dostępnych w przeglądarce, komercyjnych API chmurowych oraz modeli opartych na uczeniu głębokim, takich jak Whisper.

Hipoteza badawcza zakłada, że modele oparte na uczeniu głębokim, w szczególności rozwiązania takie jak OpenAI Whisper, zapewniają wyższą dokładność transkrypcji niż komercyjne API chmurowe, jednak mogą charakteryzować się dłuższym czasem przetwarzania danych (latency) przy przesyłaniu całych sekwencji audio. Zakres pracy obejmuje zatem nie tylko analizę jakościową wyników (WER, CER), ale również ocenę wydajności i praktycznej przydatności badanych technologii w kontekście implementacji w architekturze webowej, opartej na asynchronicznym przetwarzaniu nagrań.
\subsection{Metodyka badań i narzędzia}
Metodyka wstępnie:


\subsection{Struktura pracy}
